\subsection{可测集与可测函数}

定义 7. --- 设 f 是定义在 T 上、以拓扑空间 F（Hausdorff 与否均可）为值的函数。则称 f 对测度 \(\mu\) 可测（或称 \(\mu\) -可测），如果对 T 的每个紧子集 K，函数 \(f_{K}\) 是 \(\mu_{K}\) -可测的。

这等于说，对每个紧集 K，存在把 K 分解为一个 \(\mu_{K}\) -可忽略集 N 与一列紧集 \((K_{n})\) 的划分，使 f 到每个 \(K_{n}\) 的限制连续。由于说 N 是 \(\mu_{K}\) -可忽略的与说它局部 \(\mu\) -可忽略是等价的 (No.~4)，可见 f 是 \(\mu\) -可测的当且仅当对每个紧集 K，存在把 K 分解为一个局部 \(\mu\) -可忽略集 N 与一列紧集 \((K_{n})\) 的划分，使对一切 n，\(f_{K_{n}}\) 连续。~这一定义与 Ch. IV, §5, No.~1 的 Def. 1 相同，因而当 T 局部紧时重新得到通常的可测函数概念。

若 T 的子集 A 的特征函数可测，则称 A 可测。当 A 是 \(\mu\) -可测的且 \(\mu^{\bullet}(A) < +\infty\) 时，这个数简记为 \(\mu(A)\) ，称为 A 的测度。类似地，\(\mu(f)\) 表示 \(\mu^{\bullet}(f)\) ，只要 f 是 \(\geqslant 0\) 、\(\mu\) -可测且 \(\mu^{\bullet}(f) < +\infty\) 的函数。

若 \(\theta\) 是 \(\mathbf{T}\) 上的一个复测度，则称函数 \(f\) （相应地 \(\mathbf{T}\) 的子集）是 \(\theta\) -可测的，如果它对正测度 \(|\theta|\) 可测。下面的结果可以推广到复测度。

例. --- 设 L 是 T 的紧子集，\(\lambda\) 是 L 上的一个测度，\(\mu\) 是由 \(\lambda\) 定义的 T 上的测度 (No.~3, 例 2)。定义在 T 上的函数 f 是 \(\mu\) -可测的当且仅当 \(f_{L}\) 是 \(\lambda\) -可测的。因为，这一条件显然必要。反之，若它满足，则存在把 L 分解为一个 \(\lambda\) -可忽略集 N 与一列紧集 \((L_{n})\) 的划分，使对一切 n，\(f_{L_{n}}\) 连续。~若 K 是 T 的紧子集，则集合 \(K - \bigcup_{n}(K \cap L_{n})\)
与 L 的交是 \(\lambda\) -可忽略的，于是由 No.~3 的公式 (3)，该集是 \(\mu\) -可忽略的，且对一切 n，f 到 \(K \cap L_{n}\) 的限制连续。

Def. 7 使我们能够把关于可测函数的若干结果不加新证明地推广到空间非局部紧的情形。下面是其中一些今后将直接引用的结果：T 的开集与闭集都是 \(\mu\) -可测的；\(\mu\) -可测集构成一个部族 (Ch. IV, §5, No.~4, Th. 2 的 Cor. 2)，它包含 T 的 Borel 集 (loc. cit., Cor. 3) 以及 Souslin 集 (Ch. IV, §5, No.~1, Prop. 3 的 Cor. 2) \(^{(1)}\) 。数值函数的通常代数运算保持可测性 (Ch. IV, §5, No.~3)，可数的取极限运算亦然 (loc. cit., No.~4, Th. 2 与 Cor. 1)。下列性质值得更明确地指出：

命题 4. --- 设 \(f\) 是一个正函数，\((g_n)_{n \geqslant 1}\) 是一列 \(\mu\) -可测的正函数（定义在 \(T\) 上）。令 \(g = \sum_{n \geqslant 1} g_n\) ，则有

\[
\mu^ {\bullet} (f g) = \sum_ {n \geqslant 1} \mu^ {\bullet} (f g _ {n}).\tag{4}
\]

令 \(h_n = \sum_{i=1}^{n} g_i\) （对所有 \(n \geqslant 1\)）。对每个紧子集 \(K\) ⊂ \(T\)，

\[
\mu_ {\mathrm{K}} ^ {\bullet} \big ((f h _ {n}) _ {\mathrm{K}} \big) = \sum_ {i = 1} ^ {n} \mu_ {\mathrm{K}} ^ {\bullet} \big ((f g _ {i}) _ {\mathrm{K}} \big)
\]

（由应用于紧空间 K 的 Ch. V, §1, No.~1 的 Prop. 2）。关于 T 的紧子集的递增定向集取极限，即得

\[
\mu^ {\bullet} (f h _ {n}) = \sum_ {i = 1} ^ {n} \mu^ {\bullet} (f g _ {i}).
\]

而 \(fg\) 是递增序列 \((fh_n)_{n\geqslant 1}\) 的极限，故 \(\mu^{\bullet}(fg) = \lim_{n\to \infty}\mu^{\bullet}(fh_n)\) ；于是前面的公式立即蕴含 (4)。

推论. --- 设 \((\mathbf{A}_{n})\) 是两两不相交的可测子集序列，其并为 A。对 T 的每个子集 B，

\[
\mu^ {\bullet} (\mathrm{A} \cap \mathrm{B}) = \sum_ {n} \mu^ {\bullet} (\mathrm{A} _ {n} \cap \mathrm{B})
\]

特别地，

\[
\mu^ {\bullet} (\mathrm{A}) = \sum_ {n} \mu^ {\bullet} (\mathrm{A} _ {n}).
\]

在如上推广到 Hausdorff 空间的可测函数或可测集的性质中，我们还引用 Ch. IV, §5, No.~8 的 Prop. 12（紧集的 \(\mu\) -稠密族）。于是，以拓扑空间（Hausdorff 与否均可）为值的函数 f 是 \(\mu\) -可测的当且仅当 T 中使 \(f_{K}\) 连续的紧子集 K 所成的集合是 \(\mu\) -稠密的 (loc. cit., No.~10, Prop. 15)。

